\section{Finite three-dimensional driven hotspot packing}
\label{hs:three-section}

This section constructs finite inclusions with conical \emph{point} tips.
Small positive contrast supplies a convenient uniform scattering inverse;
it does not remove the point singularity.  The material bounds and frequency
remain fixed.  The modal question is separated at the end.

The two-dimensional construction suggests a mechanism, but truncating its
invariant direction would introduce new boundaries and radiation.  Here
each sharp reference point instead has two opposite dielectric cones
meeting at their tips.  The resulting high-field region is localized in
all three spatial directions.  We again first prove excitation of the
sharp points by one physical illumination and then round the tips.

\begin{theorem}[Finite three-dimensional driven point localization]
\label{hs:three-packing}
Fix $k>0$ and an upper bound $\varepsilon_+>1$.  There is a number
$\varepsilon_h\in(1,\varepsilon_+]$ that can be kept fixed in the
following construction.  Given any nonempty open patch $\mathcal P$ in
a plane, any finite $N$, and any $\ell,M>0$, there exist a compactly
supported binary dielectric $\varepsilon_r\in\{1,\varepsilon_h\}$
with $2N$ disjoint smooth inclusions and one unit-incoming-power
Maxwell illumination at frequency $k$ such that there are $N$
pairwise disjoint three-dimensional balls $Q_j$, of diameter less
than $\ell$, satisfying
\begin{equation}
 \max_{\overline Q_j,\,\text{one-sided}}|E|>M,
 \qquad
 \sup_{\text{a collar of }\partial Q_j,\,\text{one-sided}}|E|
 <\tfrac12\max_{\overline Q_j,\,\text{one-sided}}|E|.
 \label{hs:three-peak}
\end{equation}
The centers of the balls lie in $\mathcal P$, and the orthogonal
projections of the balls and of the material support are contained
in $\mathcal P$.  In particular the selected superlevel components
are localized in all three spatial dimensions.  The maximizing
points after rounding need not lie exactly in the original plane.
There is no invariant infinite direction, no internal impressed
source, and no conducting boundary in the physical device.
\end{theorem}

The prescribed patch is therefore a projection constraint, while the
localized neighborhoods themselves are three-dimensional.  In particular,
the theorem concerns finite point neighborhoods rather than invariant ridges.
As before, no minimum energy fraction in an individual peak is required.

The proof follows the same logic as the cross-sectional case, but each
step is a three-dimensional Maxwell statement.  A spherical angular
problem supplies a conical exponent below one; a coefficient functional
detects whether the corresponding electric singularity is present.  A
scattering inverse and total-field approximation provide an illumination
with a nonzero coefficient at every chosen point.  Finally, convergence
under rounding retains high values inside each ball, and bounds on its
collar separate the resulting maximum from the boundary.  The small
contrast is selected once within the prescribed material bound; the
number and size of the peaks may then be chosen as stated in the theorem.

\subsection{A genuine conical point exponent}
\label{hs:three-angular-subsection}

The two caps in \eqref{hs:three-caps} describe the directions occupied
by the opposing cones.  The angular eigenfunction determines the shape of
a potential $r^\lambda\phi$ near their common tip.  Because
$0<\lambda<1$, its gradient is unbounded at the tip but has finite local
energy.  This electrostatic calculation identifies the candidate
singularity; the following subsection is needed to establish its role in
the finite-frequency Maxwell field.

\begin{lemma}[Strict positive contrast gives a conical exponent below one]
\label{hs:three-angular}
Choose $t\in(1/\sqrt3,1)$ and set, on the unit sphere,
\begin{equation}
 q(\widehat x)=1+d\,\mathbf1_{\{|\widehat x\cdot e_3|>t\}},
 \qquad d>0.
 \label{hs:three-caps}
\end{equation}
There are a real $\phi\in H^1(S^2)$ and $\lambda\in(0,1)$ satisfying
\begin{equation}
 -\operatorname{div}_{S^2}(q\nabla_{S^2}\phi)
   =\nu q\phi,\qquad
 \nu=\lambda(\lambda+1),\qquad
 \int_{S^2}q\phi^2=1,\quad \int_{S^2}q\phi=0.
 \label{hs:three-angular-eigenproblem}
\end{equation}
The electrostatic field $E_0=\nabla(r^\lambda\phi)$ is curl free,
satisfies $\operatorname{div}(qE_0)=0$ weakly in a ball, and obeys
\begin{equation}
 \int_{B_R}q|E_0|^2\,dx=\lambda R^{2\lambda+1}.
 \label{hs:three-static-energy}
\end{equation}
\end{lemma}

\begin{proof}
The weighted Rayleigh quotient on the subspace orthogonal to constants
has a positive minimum by the compact embedding
$H^1(S^2)\hookrightarrow L^2(S^2)$ and the sphere Poincare inequality.
Use the trial function $z=\widehat x\cdot e_3$, whose weighted mean
vanishes by reflection symmetry.  Since
$|\nabla_{S^2}z|^2=1-z^2$,
\begin{equation}
 \int_{S^2}q(|\nabla_{S^2}z|^2-2z^2)
 =d\int_{|z|>t}(1-3z^2)<0.
 \label{hs:three-trial}
\end{equation}
Explicitly, the trial quotient is
\begin{equation}
 \mathcal Q(z)=\frac{2+d(2-3t+t^3)}{1+d(1-t^3)},\qquad
 \mathcal Q(z)-2=-\frac{3dt(1-t^2)}{1+d(1-t^3)}<0.
 \label{hs:three-trial-explicit}
\end{equation}
Thus the first positive eigenvalue $\nu$ is less than two; its
positive indicial root satisfies $0<\lambda<1$.  Separation in
spherical coordinates proves the weak electrostatic equation.
The interface conditions are exactly those in the angular weak
problem.  The flux through a small sphere is of order
$r^{\lambda+1}$, so there is no point source at the origin.
Finally $\int q|\nabla_{S^2}\phi|^2=\nu$ gives
$(\lambda^2+\nu)\int_0^Rr^{2\lambda}\,dr
=\lambda R^{2\lambda+1}$.
\end{proof}

\subsection{An exact Maxwell coefficient functional at the conical point}
\label{hs:three-functional-subsection}

At positive frequency the Maxwell field is not simply the electrostatic
field just described.  The decomposition below separates a gradient
part, which can contain the conical singularity, from a more regular
remainder.  Projecting the gradient part onto the angular eigenfunction
then defines a coefficient that can be tested for the actual total field.
This is the bridge from the local model to a physical scattering solution.

We use one standard regularity theorem in a precise local form.  On a
ball there is a bounded linear regularized Poincare lifting
\begin{equation}
 \mathcal R:H^1(B)^3\cap\ker\operatorname{div}\longrightarrow H^2(B)^3,
 \qquad \operatorname{curl}\mathcal R f=f.
 \label{hs:three-lifting}
\end{equation}
This is the regularized curl-potential operator of
Costabel and McIntosh~\cite{costabelMcIntosh2010}; their homotopy
identity and order-minus-one mapping theorem apply to a ball.
Only this established local Sobolev mapping is imported.  It is not
a scattering, corner-excitation, or rounding theorem.

Use normalized units $\operatorname{curl}E=ikH$,
$\operatorname{curl}H=-ikqE$ in the local arguments; the magnetic
field here is the physical magnetic field multiplied by $Z_0$.

\begin{lemma}[Local regular decomposition with controlled remainder]
\label{hs:three-decomposition}
Suppose $E,H$ solve the source-free Maxwell system with
$1\leq q\leq\varepsilon_h$ on $B_{3R}$ and have finite local energy.
On $B_R$ one can write
\begin{equation}
 E=\nabla v+W,\qquad v\in H^1(B_R),\quad W\in H^2(B_R)^3,
 \quad
 \|v\|_{H^1}+\|W\|_{H^2}\leq C\|E\|_{L^2(B_{3R})},
 \label{hs:three-decomposition-bound}
\end{equation}
after fixing the additive constant of $v$.  The maps can be chosen
linear.  The constant depends on the displayed radii, frequency and
material bounds, not on derivatives of $q$.
\end{lemma}

\begin{proof}
Testing $\operatorname{curl}\operatorname{curl}E=k^2qE$ against
$\chi^2E$, for a compactly supported cutoff, and applying Young's
inequality gives
$\|\operatorname{curl}E\|_{L^2(B_{5R/2})}
\leq C\|E\|_{L^2(B_{3R})}$.
Since $\operatorname{div}H=0$ and
$\operatorname{curl}H=-ikqE$, the Fourier curl/div identity applied
to a cutoff of $H$ gives
$\|H\|_{H^1(B_{2R})}\leq C\|E\|_{L^2(B_{3R})}$.
This is the same constant-permeability interior mechanism as
Lemma~\ref{ch:magneticHone}.  Apply \eqref{hs:three-lifting} on
$B_{2R}$ to $\operatorname{curl}E=ikH$, which is divergence free,
and restrict the resulting $W$ to $B_R$.
The curl-free difference $E-W$ is a gradient on the simply connected
ball.  Its potential with zero mean satisfies the stated bound by
the ordinary Poincare inequality.  All operations are bounded and
linear.
\end{proof}

The local decomposition separates the singular potential from a bounded
remainder. The coefficient below is defined by an angular projection;
its independence of the chosen lifting makes it a property of the field,
rather than of the decomposition.

\begin{lemma}[Intrinsic conical singular coefficient]
\label{hs:three-coefficient}
In the conical material of Lemma~\ref{hs:three-angular}, every local
Maxwell solution as above has a linear coefficient $c(E)$, continuous
for the $L^2(B_{3R})$ topology on its solution space.  If $c(E)\ne0$,
then $|E|$ is unbounded at ordinary points approaching the conical
vertex from the homogeneous regions.  Its local energy remains finite.
\end{lemma}

\begin{proof}
Use $E=\nabla v+W$ from Lemma~\ref{hs:three-decomposition}.  Since
$H^2(B_R)\hookrightarrow C^{0,\gamma}(\overline B_R)$ for
$0<\gamma<1/2$, $W$ is bounded.  Define
\begin{equation}
 U(r)=\int_{S^2}qv(r,\widehat x)\phi,\qquad
 A(r)=\int_{S^2}qW_r(r,\widehat x)\phi,\qquad
 B(r)=\int_{S^2}qW_t(r,\widehat x)\cdot\nabla_{S^2}\phi.
 \label{hs:three-projections}
\end{equation}
The subscripts denote radial and tangential components.  Projecting
$\operatorname{div}(q(\nabla v+W))=0$ gives
\begin{equation}
 (r^2U')'-\lambda(\lambda+1)U=-(r^2A)'+rB.
 \label{hs:three-radial}
\end{equation}
The exact solution with finite $H^1$ energy is
\begin{align}
 U(r)&=c(E)r^\lambda+p(r),\notag\\
 p(r)&=\frac{1}{2\lambda+1}\left[
 r^\lambda\int_0^rs^{-\lambda}\{B(s)-(\lambda+1)A(s)\}\,ds
 \right.\notag\\[-2pt]
 &\hspace{36mm}\left.
 -r^{-\lambda-1}\int_0^rs^{\lambda+1}\{B(s)+\lambda A(s)\}\,ds
 \right].
 \label{hs:three-exact}
\end{align}
To verify this formula, variation of constants for the two homogeneous
solutions $r^\lambda,r^{-\lambda-1}$ first gives integrals of
$-A'-2A/r+B/r$.  Integrating the $A'$ terms by parts cancels their
endpoint contributions and gives \eqref{hs:three-exact}.
The integrals converge because $A,B$ are bounded and $\lambda<1$.
They show
\begin{equation}
 |p(r)|\leq Cr\|W\|_\infty,\qquad
 |p'(r)|\leq C\|W\|_\infty.
 \label{hs:three-particular-bound}
\end{equation}
For the second inequality, differentiation produces the additional
term $-A(r)$ and two bounded integral terms.  Finite $H^1$ energy
excludes the other homogeneous term $r^{-\lambda-1}$.

Average $r^{-\lambda}(U(r)-p(r))$ on $R/2<r<R$ to define $c(E)$.
This is a bounded functional of $v$ in $L^2$ and $W$ in $L^\infty$;
Lemma~\ref{hs:three-decomposition} proves the required continuity.
It is independent of the chosen lifting.  Indeed, two liftings
differ by a bounded curl-free vector field; the corresponding
potentials differ by a Lipschitz function.  Its angular projection
is $O(r)$ after subtracting its constant value, and the weighted
mean of $\phi$ is zero.  Such a difference has coefficient zero.

Angular integration by parts gives
\begin{equation}
 \sqrt{\nu}|U(r)|\leq
 \left(\int_{S^2}q|\nabla_{S^2}v(r,\widehat x)|^2\right)^{1/2}.
 \label{hs:three-gradient}
\end{equation}
If $c(E)\ne0$, then $|U(r)|\geq |c(E)|r^\lambda/2$ at all
sufficiently small admissible radii, since $p=O(r)$.
Consequently the essential supremum of $|\nabla v|$ on these spheres
is at least a positive constant times $r^{\lambda-1}$.
Subtracting the bounded $W$ cannot remove the divergence of $|E|$.
Inside each constant-coefficient region the Maxwell field is analytic;
the large values can be selected away from interfaces.  Its energy
is finite by the starting local $L^2$ assumption.
\end{proof}

The local coefficient alone does not determine whether a prescribed
illumination excites the singularity. We next construct a local
finite-frequency solution with a nonzero coefficient; the subsequent
Runge argument supplies a physical illumination.

\begin{lemma}[Local positive-frequency excitation of the point coefficient]
\label{hs:three-local-excitation}
For the conical coefficient $q$, at every fixed $k>0$ there are local
source-free Maxwell solutions on a sufficiently small ball with
$c(E)\ne0$.
\end{lemma}

\begin{proof}
Scale a physical ball to $B_1$ and write $\kappa=kR$.
Let $E_0=\nabla(r^\lambda\phi)$ and consider
\begin{equation}
 X_q=\{w\in H_0(\operatorname{curl},B_1):
                 \operatorname{div}(qw)=0\}.
 \label{hs:three-local-space}
\end{equation}
There is a finite constant $C_q$ with
$\|w\|_2\leq C_q\|\operatorname{curl}w\|_2$ on $X_q$.
Here is a reduction to the ordinary ball estimate, which avoids an
unproved smoothness assumption on $q$.  Decompose
$w=w_T+\nabla p$ with $p\in H_0^1$ and
$\operatorname{div}w_T=0$, $w_T\in H_0(\operatorname{curl})$.
Weighted divergence yields
$\int q\nabla p\cdot\nabla\overline\psi
=-\int(q-1)w_T\cdot\nabla\overline\psi$, so
$\|\nabla p\|_2\leq d\|w_T\|_2$.
For a ball the usual Maxwell estimate is
$\|w_T\|_2\leq C\|\operatorname{curl}w_T\|_2$.
For clarity, $H^1$ membership in this assertion is not assumed.
The order-minus-one lifting also maps $L^2$ to $H^1$; lift
$\operatorname{curl}w_T$ to $V\in H^1$.  The difference
$w_T-V=\nabla p_0$ has $\Delta p_0=-\operatorname{div}V$ and
tangential boundary gradient $-V_t$.  The latter is in
$H^{1/2}(S^2)$, so the boundary trace of $p_0$, modulo a constant,
is in $H^{3/2}(S^2)$.  This follows by the scalar sphere Laplacian
estimate applied to its tangential gradient.  Dirichlet regularity
on the smooth ball gives $p_0\in H^2$, hence $w_T\in H^1$.
One can now use the Gaffney identity for tangential-zero fields,
\begin{equation}
 \|\operatorname{curl}w_T\|_2^2+\|\operatorname{div}w_T\|_2^2
 =\|\nabla w_T\|_2^2+2\|w_T\cdot n\|_{L^2(S^2)}^2,
 \label{hs:three-gaffney}
\end{equation}
initially for smooth fields and then by trace-preserving approximation
on the smooth ball.  Poincare's inequality with the tangential trace
condition excludes constant vectors and bounds the $L^2$ norm.
Thus $C_q\leq(1+d)C$ is sufficient.

Solve on $X_q$
\begin{equation}
 \int\operatorname{curl}w\cdot\operatorname{curl}\overline z
       -\kappa^2q w\cdot\overline z
 =\kappa^2\int qE_0\cdot\overline z.
 \label{hs:three-local-variational}
\end{equation}
For $\kappa^2\varepsilon_h C_q^2<1/2$ this is coercive and gives
$\|w\|_{H(\operatorname{curl})}\leq C\kappa^2\|E_0\|_2$.
It is also the equation for every test in $H_0(\operatorname{curl})$:
the weighted Helmholtz decomposition splits such a test into an
$X_q$ part and a gradient from $H_0^1$, and both sides vanish on
the gradient part.  Hence $E=E_0+w$ satisfies
$\operatorname{curl}\operatorname{curl}E=\kappa^2qE$ and
$H=(i\kappa)^{-1}\operatorname{curl}E$ completes the exact Maxwell
solution.

On an interior ball, cutoff curl/div regularity gives
\begin{equation}
 \|\operatorname{curl}w\|_{H^1(B_{1/2})}
 \leq C\left(\|\operatorname{curl}w\|_2
              +\kappa^2\|q(E_0+w)\|_2\right)=O(\kappa^2).
 \label{hs:three-local-curl-small}
\end{equation}
Apply the lifting to $\operatorname{curl}w$.  Its $W$ has $H^2$
norm $O(\kappa^2)$; the potential in $E=\nabla v+W$ differs from
$r^\lambda\phi$, up to a constant, by $O(\kappa^2)$ in $H^1$.
The averaged formula in Lemma~\ref{hs:three-coefficient} therefore
gives $c(E)=1+O(\kappa^2)$.  This proves nontriviality for small
positive $\kappa$.  The local tangential boundary data used to
construct this solution are a mathematical matching device.  The
physical construction below uses global incoming waves and has no
PEC boundary.
\end{proof}

\subsection{A uniform finite-device Maxwell inverse at fixed small contrast}
\label{hs:three-volume-subsection}

The following inverse controls the complete response of the finite
dielectric, including multiple scattering among the inclusions.  Small
contrast makes the volume operator invertible uniformly over the family
used for rounding.  It is a device for controlling the global problem;
the local conical exponent remains below one at every positive contrast.

\begin{lemma}[Uniform volume-integral inverse]
\label{hs:three-volume}
Fix a ball $B_L$ and $k>0$.  There is $C_{kL}<\infty$ such that
the electric volume operator
\begin{equation}
 \mathcal K_k f=(k^2I+\nabla\nabla)\Phi_k*f,
 \qquad \Phi_k(x)=\frac{e^{ik|x|}}{4\pi|x|},
 \label{hs:three-volume-operator}
\end{equation}
restricted to $B_L$, with $f$ zero extended, has $L^2$ norm at most
$C_{kL}$.  For $0<d<(2C_{kL})^{-1}$ and any measurable
$0\leq\chi\leq d$ supported in $B_L$,
\begin{equation}
 E=E^{\rm in}+\mathcal K_k(\chi E)
 \label{hs:three-volume-equation}
\end{equation}
has a unique solution and
$\|(I-\mathcal K_k\chi)^{-1}\|\leq2$.
It is the unique outgoing Maxwell scattering field for
$\varepsilon_r=1+\chi$.  An outgoing electric source problem with
compact $L^2$ forcing is likewise well posed.
\end{lemma}

\begin{proof}
Split
$\mathcal K_k=\nabla\nabla\Phi_0*
 +[k^2I\Phi_k+\nabla\nabla(\Phi_k-\Phi_0)]*$.
The first operator has Fourier multiplier
$-\xi\xi^T/|\xi|^2$, of norm at most one on all-space $L^2$;
this includes its distributional diagonal term.  The bracketed
matrix kernel is locally integrable: the expansion
$\Phi_k-\Phi_0=ik/(4\pi)-k^2|x|/(8\pi)+O(k^3|x|^2)$
shows a Hessian of order $1/|x|$.  Away from zero it is smooth.
Young's inequality after cutting the kernel off at $|x|=2L$
therefore bounds the remainder by its finite matrix $L^1$ norm.
Rescaling $x$ by $L$ shows that the resulting constant depends only
on $kL$.  This proves the operator estimate.

The Neumann series proves the inverse bound, uniformly over all
the admitted geometries.  The dyadic kernel
$k^{-2}\mathcal K_k$ is the outgoing fundamental solution of
$\operatorname{curl}\operatorname{curl}-k^2$.  Consequently
\eqref{hs:three-volume-equation} gives the Maxwell equations and
radiation condition in distributions.  Conversely, subtracting
that representation from an outgoing solution gives an entire
outgoing vacuum solution, which is zero by the spherical-wave
expansion.  This proves uniqueness and equivalence, including for
discontinuous $\chi$.  An electric forcing $f\in L^2_c$ replaces
$E^{\rm in}$ by $k^{-2}\mathcal K_k f$ on a ball enclosing its
support; the same inverse applies when that support is in $B_L$.
Only this latter case is needed below.
\end{proof}

For the theorem choose $L=1/k$ and fix once and for all
\begin{equation}
 0<d<\min\{\varepsilon_+-1,(2C_1)^{-1}\},
 \qquad \varepsilon_h=1+d.
 \label{hs:three-fixed-contrast}
\end{equation}
Every finite collection of tips and inclusions can be placed in a
much smaller ball whose projected closure lies in $\mathcal P$.
The containing comparison ball of radius $1/k$ is only used for
the inverse; it is not extra material.

\subsection{Physical illumination, total-field density, and rounding}
\label{hs:three-global-subsection}

The incoming field is a superposition of plane waves with a square
integrable angular density.  Its incoming power is finite and can be
normalized to one.  The density result below concerns the resulting
\emph{total} fields in the actual dielectric, not just the incident waves.
It connects the local conical solutions to one globally admissible source.

We now restore physical units: $E$ is the physical electric field
and $H$ is the physical magnetic field, without the local $Z_0$
rescaling used above.  For a smooth tangential density
$g:S^2\to\mathbb C^3$,
$d\cdot g(d)=0$, use the physical electric incident field
\begin{equation}
 E_g^{\rm in}(x)=\frac{k\sqrt{Z_0}}{\sqrt2\,\pi}
                 \int_{S^2}g(d)e^{ikd\cdot x}\,dS(d),
 \qquad
 H_g^{\rm in}(x)=\frac{k}{\sqrt{2Z_0}\,\pi}
                 \int_{S^2}d\times g(d)e^{ikd\cdot x}\,dS(d).
 \label{hs:three-illumination}
\end{equation}
Stationary phase on the sphere gives an incoming electric amplitude
$\sqrt{2Z_0}\,g(-\widehat x)/r$, up to a phase of unit modulus.
The radial incoming flux is $|E^{\rm in}_{\rm incoming}|^2/(2Z_0)$,
so the total incoming power is exactly $\|g\|_{L^2(S^2)}^2$.
This is finite total power, rather than power per invariant length.

\begin{lemma}[Total-field Maxwell Runge approximation for these devices]
\label{hs:three-runge}
Let the small-contrast coefficient be piecewise constant, with finitely
many smooth interface pieces and finitely many conical vertices and
joining curves.  Let $Q$ be a closed ball, or a finite disjoint union
of closed balls, compactly contained in the comparison ball $B_L$
and with connected complement.  A source-free Maxwell
electric field defined on a neighborhood of $Q$ can be approximated
in $L^2(Q)^3$ by total electric fields from
\eqref{hs:three-illumination}.
\end{lemma}

\begin{proof}
Take $f\in L^2(Q)^3$ annihilating all such total fields under the
complex bilinear pairing, and solve
$(\operatorname{curl}\operatorname{curl}-k^2\varepsilon_r)W=f$
outgoing by Lemma~\ref{hs:three-volume}.  Orthogonality to every
smooth tangential density says
$\int_Q f\cdot E_{d,p}=0$ for all incident plane waves
$p e^{ikd\cdot x}$ with $p\cdot d=0$.  Continuity in $d,p$
justifies passing from integrated illuminations to individual
plane waves.  Maxwell Green's identity, without complex conjugation,
equates this integral to a nonzero constant times
$p\cdot W_\infty(-d)$.  One may obtain the identity directly by
inserting the dyadic outgoing representation; the pairing of two
outgoing solutions has zero boundary Wronskian.  Since the electric
far field is tangential, testing all $p\perp d$ proves
$W_\infty=0$.  The vacuum vector spherical-wave expansion then
gives $W=0$ outside a sufficiently large sphere.

The exterior of $Q$ is connected.  Unique continuation here is
elementary within each material: the electric and magnetic fields
satisfy constant-coefficient Helmholtz equations.  Zero extends
across an open smooth interface because both tangential electric
and magnetic traces are zero; extending by zero gives a weak
constant-coefficient Maxwell solution on a ball crossing that
interface, hence an analytic one.  Paths avoid the finitely many
vertices and joining curves and cross only open interface pieces.
Thus $W=0$ outside $Q$.
Both $W$ and $\operatorname{curl}W$ are in local
$H(\operatorname{curl})$, and have zero tangential traces on
$\partial Q$.  For every local homogeneous solution $V$, Green's
identity consequently gives
$\int_Q f\cdot V=\int_QW\cdot
(\operatorname{curl}\operatorname{curl}-k^2\varepsilon_r)V=0$.
The annihilator characterization of the $L^2$ closure proves density.
This is density of total fields in the actual material, not an
assumed free-space approximation inside it.
\end{proof}

The rounding estimate has two distinct roles: $L^2$ convergence transfers
large values from the sharp reference field, while the collar bound confines
the resulting peaks. The latter depends on regularity estimates in fixed
smooth-interface charts away from the shrinking tips; convergence in
$L^2$ alone does not supply a pointwise bound on the collars.

\begin{lemma}[Finite-device rounding continuity]
\label{hs:three-rounding}
Let $\chi_\rho,\chi$ be supported in the common comparison ball,
$0\leq\chi_\rho,\chi\leq d$, with $\chi_\rho\to\chi$ almost
everywhere.  For one fixed incident field, their total electric
fields satisfy $E_\rho\to E$ in $L^2(B_L)^3$.
For the smooth rounded conical family used here, suppose in addition
that the coefficients agree exactly with one fixed smooth reference
body on a neighborhood of each prescribed sphere collar, for every
sufficiently small rounding parameter.  The collars avoid all vertices
and joining curves.  Then the continuous electric fields on both closed
material sides have one common finite bound on those collars.  This
includes collars intersecting smooth conical interfaces.
\end{lemma}

\begin{proof}
Subtract the volume equations:
\begin{equation}
 E_\rho-E=(I-\mathcal K_k\chi_\rho)^{-1}
                  \mathcal K_k((\chi_\rho-\chi)E).
 \label{hs:three-rounding-identity}
\end{equation}
Dominated convergence makes the last multiplication tend to zero
in $L^2$, and Lemma~\ref{hs:three-volume} bounds both operators
uniformly.  This proves the first claim.

For the collar bound, write the actual total field as
$E_\rho=\nabla v_\rho+W_\rho$, retaining the same fixed smooth
incident field.  The convergent volume unknowns have bounded source
norms; the compact scalar Green operators give bounded local scalar
potential and gradient norms, and bounded $H^2$ correction norms.
Choose a fixed smooth rounded reference body that agrees with every
later body near the collar.  Its regular interfaces supply fixed
signed graph charts; homogeneous neighborhoods complete the cover.

In each chart the actual scalar transmission equation is localized
with nested smooth cutoffs.  A regional difference-quotient energy
estimate first derives the mixed derivative entering the differentiated
forcing.  Thus no second derivative of the original potential on an
outer cutoff annulus is assumed.  Two tangential iterations and the
normal-flux equation then give continuous tangential and flux
representatives with bounds depending only on the original scalar and
forcing norms and the fixed chart data.  Reconstructing the physical
gradient gives uniformly bounded actual electric branches on both
closed material sides.  A finite cover of the compact sphere gives
one bound on an open collar, chosen before the rounding index.
The same local regularity argument, without a uniform constant, gives
continuous one-sided branches for every fixed smooth rounded device.
No estimate uniform at a shrinking tip is used.
\end{proof}

\begin{proof}[Proof of Theorem~\ref{hs:three-packing}]
Select distinct centers $p_j$ in the planar patch and sufficiently
small disjoint balls about them.  Put two opposite dielectric cones
at each center, with opening angle from \eqref{hs:three-caps}.
They agree exactly with the homogeneous conical geometry on a
smaller ball, and are smoothly capped farther away, all within
the selected neighborhood.  Their closures touch only at the
center in this sharp reference geometry.  The projections of all
these neighborhoods lie in $\mathcal P$, and their total support
fits in a ball of radius $1/k$.

For each selected vertex, Lemma~\ref{hs:three-local-excitation}
constructs a local finite-frequency field with nonzero coefficient.
Lemma~\ref{hs:three-runge} approximates it on a slightly larger
ball than that used to measure the coefficient.
Lemma~\ref{hs:three-coefficient} then shows that the coefficient
$c_j(E_g)$ is a nonzero continuous linear functional of the smooth
tangential incident density $g$.  Choose one density outside the
finite union of its kernels, exactly as in
Proposition~\ref{hs:simultaneous}, and divide by its $L^2(S^2)$ norm.
All coefficients remain nonzero, and the incoming power is one.

Fix disjoint observation balls $Q_j$ of diameter less than $\ell$
inside the exact conical regions, and round all tips inward while
leaving their collars unchanged.  The result has $2N$ disjoint smooth
inclusions and the same two material values.  The changed sets shrink
to the finite vertex set.  Lemma~\ref{hs:three-rounding} gives a common
bound $B$ for both actual material-side traces on all fixed collars,
uniformly for sufficiently small rounding parameters.

Choose $S>\max\{M,2B\}$.  Each nonzero conical coefficient supplies
a positive-volume superlevel $|E|>2S$ inside its observation ball.
Local $L^2$ convergence preserves a positive-volume superlevel
$|E_\rho|>S$ in every ball for all sufficiently small $\rho$.
The continuous one-sided branches attain their joint maximum on the
compact ball.  Each maximum exceeds $S$ and twice the collar bound.
The collar is therefore a strict low barrier for its half-amplitude
superlevel set.  Its full connected component is compact and stays
strictly inside the observation ball, proving \eqref{hs:three-peak}.
All choices of illumination precede rounding; the same normalized
input is used for every rounded field.
\end{proof}

\begin{corollary}[Finite-region volume-to-peak ratios in three dimensions]
\label{hs:three-regional-volume}
Along the rounded sequence for one fixed sharp configuration and
one normalized input, electromagnetic energy in a fixed bounded
region stays bounded, while the electric maximum in every selected
ball tends to infinity.  Each resulting finite-region energy-to-peak
ratio tends to zero.  This is a driven local-energy statement, not
an all-space modal normalization.
\end{corollary}

\begin{proof}
The electric $L^2$ norms are uniformly bounded by the volume inverse.
The magnetic fields are bounded on fixed bounded regions by the
Maxwell equations, the local curl estimate, and exterior vacuum
regularity; a slightly larger comparison ball may be used for the
cutoffs.  For every finite amplitude threshold, the sharp field has a
positive-volume superlevel in every selected ball.  Local $L^2$
convergence preserves a smaller superlevel for every sufficiently late
rounded device.  The attained one-sided maxima therefore diverge
along the same family.  The material factors are bounded above and below.
\end{proof}

\subsection{What an exact mode can mean in a finite dielectric device}
\label{hs:three-modal-subsection}

\begin{proposition}[Absence of real-frequency finite-energy modes]
\label{hs:three-no-real-mode}
Let the real positive scalar permittivity be constant on finitely
many material pieces with piecewise smooth interfaces, and equal
one outside a bounded set.  At a fixed real $\omega>0$, a source-free
Maxwell field with finite all-space electromagnetic energy is zero.
The same conclusion holds for an outgoing source-free field,
without a finite-energy assumption.
\end{proposition}

\begin{proof}
In the exterior vacuum every Cartesian component satisfies the
scalar Helmholtz equation.  A spherical harmonic expansion of an
$L^2$ exterior solution has radial terms proportional to incoming
and outgoing spherical Hankel functions.  Each nonzero radial
combination has nondecaying mean radial $L^2$ mass, because its
leading amplitude is a nonzero combination of $e^{ikr}$ and
$e^{-ikr}$ divided by $r$.  It cannot be square integrable on
$(R,\infty)$ with volume weight $r^2$.  Orthogonality of spherical
harmonics excludes cancellation between different angular terms.
Thus a finite-energy field is zero in the exterior.
For an outgoing source-free field, the real part of the complex
Poynting identity instead sets its nonnegative far-field radiation
power to zero, giving the same exterior conclusion.
Tangential electric and magnetic continuity then propagates zero
across each open smooth material interface by the constant-coefficient
extension argument used in Lemma~\ref{hs:three-runge}.
Paths can avoid the finite lower-dimensional interface singular sets.
Hence the field is zero everywhere.
\end{proof}

This proposition diagnoses an inappropriate literal interpretation
of the modal numerator; it does \emph{not} refute the tutorial's
high-$Q$ or quasinormal-mode interpretation.  A finite radiating
resonator has complex-frequency poles.  Their outgoing fields grow
at spatial infinity and require a specified regularized mode norm,
or a justified high-$Q$ approximation, instead of the ordinary
positive all-space energy integral.  The tutorial explicitly notes
open-cavity generalizations of its mode-volume definition.

\begin{remark}[The driven proof and the separate modal construction]
\label{hs:three-modal-gap}
Theorem~\ref{hs:three-packing} establishes the finite-device \emph{driven}
point-localization statement given above.  It does not turn the chosen
illumination into a resonator mode.  A possible modal route starts
with a simple isolated dielectric resonance, inserts finitely many
small conical inclusions where its electric field is nonzero, and
then rounds their tips.  To complete that route one must prove
that a single continued resonance has nonzero conical coefficients
at every selected vertex, and that the chosen resonance normalization
stays finite and nonzero under insertion and rounding.  These are
not consequences of Runge approximation: Runge varies external
illumination, whereas a pole's eigenfield is fixed by its homogeneous
outgoing eigenproblem.  Strong local field convergence without a
coefficient topology also does not prove nonzero singular coefficients.
Sections~\ref{hs:modal-next-section} and~\ref{ms:section} establish
these conditions for the particular modal family considered there.
Eventual nonzero excitation
is proved by a strict annular pairing in
Theorem~\ref{hs:modal-coefficient-continuation}, without requiring
continuation of the intrinsic coefficient through zero contrast; a simple pole,
its genuine reciprocal normalization and its rounded continuation
are proved in Lemmas~\ref{ms:sphere-simple}--\ref{ms:rounding}.
The resulting closed modal statement is
Theorem~\ref{ms:modal-packing}.  It does not rely on varying a driven
illumination or on a divergent all-space positive energy integral.
\end{remark}
